\chapter{Rhinoplasty and Septorhinoplasty}

\section*{Introduction \& Clinical Indications}
Rhinoplasty is a sophisticated surgical procedure aimed at reshaping the external nose to address both aesthetic and functional concerns. This operation involves the careful modification of the underlying bony and cartilaginous structures of the nose, with the objective of achieving a balanced facial appearance while ensuring optimal nasal airflow. Septorhinoplasty is a specific variant that combines rhinoplasty with septoplasty, the latter focusing on correcting a deviated nasal septum, which is the partition of bone and cartilage that divides the nasal cavity. This dual approach is often necessary, as deviations in the septum can lead to both breathing difficulties and visible asymmetry of the nose. The procedure requires a deep understanding of nasal anatomy, surgical finesse, and precise technical execution to achieve stable, predictable, and aesthetically pleasing long-term results.

\begin{itemize}
  \item Nose job
  \item Nasal reshaping surgery
  \item Septoplasty (when only the septum is corrected)
  \item Functional rhinoplasty (primarily for breathing improvement)
  \item Cosmetic rhinoplasty (primarily for aesthetic enhancement)
  \item Revision rhinoplasty (for secondary procedures)
\end{itemize}

The fundamental principle of rhinoplasty is to achieve a harmonious and functional nasal form by precisely modifying the underlying osteocartilaginous framework, allowing the overlying soft tissue envelope to redrape smoothly. The rationale is twofold: to correct aesthetic disharmonies that cause patient dissatisfaction and to alleviate functional impairments, primarily nasal airway obstruction. This often involves a delicate balance, as aggressive aesthetic maneuvers can compromise breathing, and vice versa.

Technically, the procedure aims to reduce or augment specific areas of the nose, refine the nasal tip, straighten the dorsum, and correct any asymmetry. For septorhinoplasty, the core principle involves repositioning or excising portions of the deviated septal cartilage and bone to create a patent airway, often utilizing harvested septal cartilage for structural grafts to support the external nose. Key technical goals include creating a stable, well-supported nasal framework that resists collapse, establishing adequate nasal valve patency, achieving a natural-looking result that is in proportion with the rest of the face, and ensuring long-term stability of the surgical changes. The surgeon acts as an architect, sculpting the underlying structures to achieve the desired external contour and internal function.

The history of nasal reconstruction dates back to ancient India, where the surgeon Sushruta described techniques for repairing amputated noses using forehead skin flaps around 600 BC. This early form of reconstructive rhinoplasty, known as the Indian method, was later reintroduced to Europe in the 15th century. In the 16th century, Gaspare Tagliacozzi, an Italian surgeon, pioneered the use of an arm flap for nasal reconstruction, known as the Italian method.

Modern rhinoplasty, particularly for aesthetic purposes, began to emerge in the late 19th and early 20th centuries. Jacques Joseph, a German surgeon, is widely recognized as the father of modern aesthetic rhinoplasty. In the early 1900s, he developed and popularized techniques for reducing dorsal humps and refining the nasal tip through intranasal (closed) incisions, establishing many foundational principles still in use today. Throughout the mid-20th century, surgeons like Maurice Cottle and John Converse further advanced the understanding of nasal physiology and functional aspects, emphasizing the importance of the nasal airway. The late 20th century saw the increasing popularity of the "open" rhinoplasty approach, championed by surgeons such as Jack Sheen and Richard Webster, which allowed for direct visualization and more precise manipulation of the nasal framework through a small external columellar incision. More recently, the field has evolved towards "structural rhinoplasty," emphasizing cartilage grafting for long-term support, and "preservation rhinoplasty," focusing on maintaining existing structures where possible, reflecting a continuous refinement of techniques for improved predictability and natural outcomes.

Rhinoplasty and septorhinoplasty are indicated for a variety of functional and aesthetic concerns, often simultaneously.

\begin{itemize}
  \item **Nasal Airway Obstruction:** Primarily due to a deviated nasal septum, turbinate hypertrophy, or nasal valve collapse, significantly impairing breathing.
  \item **Dorsal Hump Reduction:** Correction of an overly prominent nasal bridge, often a significant aesthetic concern.
  \item **Nasal Tip Refinement:** Addressing a bulbous, drooping, or asymmetric nasal tip to achieve better definition and projection.
  \item **Crooked Nose Correction:** Straightening a nose that deviates from the midline, often a result of trauma or congenital factors.
  \item **Post-traumatic Nasal Deformity:** Reconstruction of the nose following fractures or other injuries that have altered its shape or function.
  \item **Congenital Nasal Deformities:** Correction of birth defects affecting nasal structure, such as those associated with cleft lip and palate.
  \item **Saddle Nose Deformity:** Augmentation of a collapsed nasal dorsum, often requiring cartilage grafting, which can result from trauma, infection, or previous surgery.
  \item **Revision Rhinoplasty:** To correct unsatisfactory outcomes or complications from a previous nasal surgery, such as persistent asymmetry, breathing issues, or over-resection.
  \item **Nasal Valve Incompetence:** Surgical stabilization of the internal or external nasal valve, often diagnosed with the Cottle maneuver, to prevent collapse during inspiration.
\end{itemize}

Septorhinoplasty holds a primary position as the definitive surgical treatment for significant nasal airway obstruction caused by anatomical deformities, particularly a deviated nasal septum or nasal valve collapse, when conservative measures like nasal sprays or breathing strips have proven ineffective. In these functional cases, it is often considered a first-line surgical intervention. Similarly, for patients seeking to address specific aesthetic concerns about their nose, rhinoplasty is the primary elective procedure to achieve their desired cosmetic outcome.

However, the procedure can also serve a secondary or reconstructive role. For instance, in cases of post-traumatic nasal deformities, congenital anomalies, or deformities resulting from previous unsuccessful nasal surgeries, rhinoplasty becomes a reconstructive or revision procedure. Revision rhinoplasty, in particular, is a secondary intervention, typically more complex due to altered anatomy, scar tissue, and compromised tissue planes from prior surgery. The decision to proceed with rhinoplasty, whether primary or secondary, is always made after thorough patient evaluation, discussion of goals, and assessment of risks and benefits, ensuring it aligns with the patient's overall health and psychological well-being.

\section*{Relevant Surgical Anatomy}
The anatomy of the nose is intricate, comprising a combination of bone, cartilage, and soft tissue, all of which are closely associated with the surrounding facial structures and vital neurovascular components. The bony framework consists of paired nasal bones at the superior aspect, which articulate with the frontal bone and the frontal processes of the maxillae laterally. The cartilaginous framework includes the paired upper lateral cartilages, which connect to the nasal bones, and the paired lower lateral cartilages (alar cartilages), which shape the nasal tip and alar rims. The nasal septum, a crucial structure, is formed by the perpendicular plate of the ethmoid bone superiorly, the vomer inferiorly, and the quadrangular septal cartilage anteriorly, which serves as a primary source for grafting material during rhinoplasty.

The soft tissue envelope consists of skin, subcutaneous fat, and a superficial musculoaponeurotic system (SMAS)-like layer that covers the underlying framework. Important muscles include the procerus, nasalis, and depressor septi nasi, which can affect nasal expression. The arterial supply is abundant, originating from both the internal and external carotid systems. The anterior and posterior ethmoidal arteries, branches of the ophthalmic artery, supply the superior nasal cavity and dorsum, while the facial artery supplies the external nose and septum through branches such as the angular, superior labial, and lateral nasal arteries. Venous drainage parallels the arterial supply, draining into the anterior facial vein and ophthalmic veins. Sensory innervation is primarily from the trigeminal nerve (V1 and V2 divisions), with the external nasal nerve (from V1) supplying the nasal tip and dorsum, and the infraorbital nerve (from V2) supplying the lateral ala. Lymphatic drainage from the external nose typically flows to the submandibular and preauricular lymph nodes. Key surgical landmarks include the radix (nasofrontal angle), nasion (deepest point of the radix), rhinion (osseocartilaginous junction), supratip break, tip defining points, columella, and alar base. The internal nasal valve, formed by the septum, caudal edge of the upper lateral cartilage, and floor of the nose, is critical for airflow regulation.

\section*{Preoperative Workup \& Preparation}
A comprehensive preoperative workup is essential for safe and successful rhinoplasty. This begins with a detailed medical history, including any prior nasal trauma or surgery, allergies, and current medications. A thorough physical examination of the internal and external nose is performed, assessing skin quality, cartilage strength, septal deviation, and nasal airway patency (e.g., Cottle maneuver). Standardized photographic documentation from multiple views is crucial for surgical planning and objective outcome assessment.

\begin{itemize}
  \item **Medical Clearance:** Patients undergo routine preoperative medical evaluation, including blood tests (complete blood count, coagulation profile, basic metabolic panel) and, if indicated by age or comorbidities, an electrocardiogram and chest X-ray to ensure fitness for general anesthesia.
  
  \item **Imaging Studies:** While not always necessary, a computed tomography (CT) scan of the paranasal sinuses may be ordered for complex septal deviations, chronic sinusitis, or to evaluate the bony architecture in revision cases. Nasal endoscopy can provide detailed visualization of the internal nasal structures.
  
  \item **Medication Adjustment:** Patients are instructed to discontinue aspirin, nonsteroidal anti-inflammatory drugs (NSAIDs), vitamin E, herbal supplements (e.g., ginkgo biloba, ginseng, garlic), and any other medications that can increase bleeding risk for at least two weeks prior to surgery.
  
  \item **Smoking Cessation:** Complete cessation of smoking for at least 4-6 weeks preoperatively is strongly advised, as smoking significantly impairs wound healing and increases the risk of complications such as skin necrosis and infection.
  
  \item **NPO Status:** Patients must adhere to strict nil per os (NPO) guidelines, typically fasting from food and liquids for 6-8 hours before surgery.
  
  \item **Prophylactic Antibiotics:** Intravenous prophylactic antibiotics are often administered prior to incision to reduce the risk of surgical site infection.
  
  \item **Informed Consent:** A detailed discussion covers the surgical plan, expected outcomes, potential risks, benefits, alternatives, and the possibility of revision surgery, ensuring the patient has realistic expectations and provides fully informed consent.
\end{itemize}

**Absolute Contraindications:**
\begin{itemize}
  \item Uncontrolled severe systemic medical conditions (e.g., unstable angina, recent myocardial infarction, severe uncontrolled diabetes, significant coagulopathy) that pose an unacceptable anesthetic or surgical risk.
  
  \item Active nasal or sinus infection, which must be treated and resolved prior to surgery.
  
  \item Unrealistic patient expectations or severe body dysmorphic disorder (BDD) that has not been adequately addressed with psychiatric intervention.
  
  \item Skeletal immaturity in very young patients, as nasal growth is not yet complete, potentially leading to unpredictable results and growth disturbances.
\end{itemize}

**Relative Contraindications \& Precautions:**
\begin{itemize}
  \item Heavy smoking, which significantly compromises wound healing and increases the risk of skin complications and infection. Patients should be strongly encouraged to quit.
  
  \item Certain autoimmune or connective tissue disorders (e.g., Wegener's granulomatosis, relapsing polychondritis) that can affect cartilage integrity and healing.
  
  \item History of keloid or hypertrophic scar formation, particularly in patients with darker skin types, as this can affect external incision healing.
  
  \item Prior extensive nasal surgery or trauma, which can lead to significant scar tissue, altered anatomy, and limited cartilage availability for grafting, making revision surgery more challenging and less predictable.
  
  \item Chronic use of medications that thin the blood, requiring careful management and temporary cessation under medical supervision.
  
  \item Psychological instability or significant psychiatric comorbidities that may impact patient compliance or satisfaction with the surgical outcome.
\end{itemize}

\section*{Surgical Technique \& Procedure}
Rhinoplasty and septorhinoplasty are typically performed under general anesthesia, though some minor procedures may utilize local anesthesia with intravenous sedation. The patient is positioned supine with the head slightly elevated. The choice between an open or closed approach depends on the surgeon's preference and the complexity of the required modifications.

\begin{itemize}
  \item **Anesthesia and Infiltration:** General anesthesia is induced. The nose and surrounding areas are infiltrated with a local anesthetic solution containing a vasoconstrictor (e.g., lidocaine with epinephrine) to provide additional pain control and minimize bleeding.
  
  \item **Incisions and Exposure:**
    \begin{itemize}
      \item **Open Approach:** A small, inverted V or W-shaped incision is made across the columella, connecting to marginal incisions inside the nostrils. The skin and soft tissue envelope are then carefully elevated from the underlying osteocartilaginous framework, providing direct, panoramic visualization.
      \item **Closed Approach:** All incisions are made inside the nostrils (e.g., intercartilaginous, infracartilaginous, transfixion incisions), and the skin is undermined through these limited access points.
    \end{itemize}
  
  \item **Septoplasty (if indicated):** Mucoperichondrial flaps are elevated from the septal cartilage and bone. Deviated portions of the septum (cartilage and/or bone) are resected, straightened, or repositioned to create a clear airway. Harvested septal cartilage is often preserved for grafting.
  
  \item **Dorsal Modification:** A dorsal hump, if present, is reduced using osteotomes, rasps, or powered instruments. The nasal bones may be narrowed through controlled osteotomies (medial and lateral) to close an "open roof" deformity created by hump reduction and to narrow a wide bony dorsum.
  
  \item **Tip Refinement:** The lower lateral cartilages are meticulously reshaped using a combination of sutures (e.g., dome-binding, interdomal, lateral crural steal sutures) and/or cartilage grafts (e.g., columellar strut, shield graft, lateral crural overlay grafts) to achieve desired projection, rotation, and definition.
  
  \item **Nasal Valve Reconstruction:** If nasal valve collapse is present, it is addressed by placing spreader grafts or spreader flaps between the upper lateral cartilages and the septum, or alar batten grafts to support the external valve.
  
  \item **Alar Base Reduction (if needed):** If the nostrils are too wide, small wedges of tissue may be excised from the alar base to narrow them.
  
  \item **Closure and Stabilization:** The skin and soft tissue envelope are redraped. Incisions are meticulously closed with fine sutures. Internal splints (e.g., silicone sheets) may be placed to stabilize the septum and prevent adhesions. An external thermoplastic or plaster splint is applied to protect the nose and maintain its new shape during the initial healing phase, often secured with surgical tape.
\end{itemize}

\section*{Complications \& Risks}
While generally safe, rhinoplasty and septorhinoplasty carry potential risks and complications, which can be broadly categorized as general surgical risks and procedure-specific risks.

**General Surgical Complications:**
\begin{itemize}
  \item **Anesthesia Risks:** Adverse reactions to anesthetic agents, including nausea, vomiting, respiratory depression, or, rarely, anaphylaxis.
  
  \item **Bleeding:** Postoperative epistaxis (nosebleed), which is usually mild but can occasionally be severe, or the formation of a septal hematoma (blood clot within the septum), which requires prompt drainage to prevent cartilage necrosis.
  
  \item **Infection:** Cellulitis of the nasal skin, abscess formation, or, rarely, osteomyelitis or chondritis.
  
  \item **Poor Wound Healing:** Delayed healing or dehiscence of incisions, particularly in smokers or those with compromised immune systems.
\end{itemize}

**Procedure-Specific Complications:**
\begin{itemize}
  \item **Unsatisfactory Aesthetic Outcome:** This is the most common concern and can include asymmetry, over-resection (e.g., saddle nose, pinched tip), under-resection (e.g., residual hump, bulbous tip), pollybeak deformity (supratip fullness), inverted V deformity, or an unnatural appearance.
  
  \item **Persistent or New Nasal Obstruction:** May result from inadequate correction of septal deviation, nasal valve collapse, turbinate hypertrophy, or excessive scarring.
  
  \item **Septal Perforation:** A hole in the nasal septum, which can cause crusting, whistling, or recurrent epistaxis, especially after aggressive septal surgery.
  
  \item **Numbness:** Temporary or, rarely, permanent numbness of the nasal tip, upper lip, or upper incisor teeth due to nerve injury.
  
  \item **Skin Complications:** Skin discoloration, telangiectasias (spider veins), or, very rarely, skin necrosis, particularly with aggressive undermining or excessive pressure from splints.
  
  \item **Anosmia or Hyposmia:** Loss or reduction of the sense of smell, usually temporary due to swelling, but rarely permanent.
  
  \item **Cerebrospinal Fluid (CSF) Leak:** An extremely rare but serious complication, typically associated with aggressive superior septal dissection near the cribriform plate.
  
  \item **Vision Changes/Blindness:** An exceedingly rare but devastating complication, usually associated with injury to the ethmoidal arteries during osteotomies, leading to retrobulbar hematoma and optic nerve compression.
  
  \item **Need for Revision Surgery:** Due to residual deformities, complications, or patient dissatisfaction, revision rates range from 5\% to 15\%.
\end{itemize}

\section*{Postoperative Care \& Recovery}
The postoperative recovery from rhinoplasty and septorhinoplasty is a gradual process, with initial healing occurring over weeks and final results emerging over many months.

Immediately after surgery, patients are monitored in the Post-Anesthesia Care Unit (PACU) for airway patency, vital signs, and pain control. Most patients are discharged home the same day. For the first 24-48 hours, patients will experience significant nasal congestion, mild to moderate pain (managed with oral analgesics), and swelling and bruising around the eyes and nose, which typically peak on day 2 or 3. Head elevation, cold compresses, and avoiding strenuous activity are crucial during this period. Internal splints or packing, if used, contribute to the feeling of congestion.

The external splint and any external sutures are usually removed around 5 to 7 days post-op. At this point, much of the bruising will have started to fade, but significant swelling, particularly of the nasal tip, will persist. Patients are advised to avoid blowing their nose, wearing glasses that rest on the bridge of the nose, and engaging in strenuous activities for at least 2-4 weeks. Light activity and return to work are generally possible within 1-2 weeks. While the majority of visible swelling resolves within the first 3-6 months, subtle swelling, especially in the nasal tip, can persist for up to a year or even longer, meaning the final aesthetic result is not fully apparent until then. Patients are counseled on the importance of patience throughout this extended healing phase.

During rhinoplasty and septorhinoplasty, the surgeon meticulously documents the intraoperative findings, which are crucial for understanding the patient's specific anatomy and the extent of correction performed. These findings typically include the nature and severity of septal deviation (e.g., C-shaped, S-shaped, caudal deviation), the size and composition of any dorsal hump (bony, cartilaginous, or mixed), the quality and strength of the nasal cartilages (upper and lower lateral cartilages), the presence of nasal valve collapse, and any asymmetries of the bony or cartilaginous framework. The amount and location of cartilage harvested for grafting (typically from the septum, but sometimes from the ear or rib) are also noted.

Pathology reports are generated for any resected tissue, although in most primary rhinoplasties, the excised bone and cartilage are sent for routine histopathological examination primarily to confirm their benign nature. This is a standard safety measure to rule out any unexpected neoplastic or inflammatory processes, such as chondroma, osteoma, or granulomatous disease, which are rare but important to exclude. The report typically confirms the presence of normal nasal cartilage and/or bone, ensuring that no underlying pathology contributed to the nasal deformity.

A normal and successful postoperative course following rhinoplasty or septorhinoplasty is characterized by a predictable progression of healing and gradual resolution of symptoms. Initial discomfort and pain are typically mild to moderate and well-controlled with oral analgesics, subsiding significantly within the first few days. Nasal congestion, often due to internal swelling and splints, gradually improves over the first week. Bruising around the eyes and nose resolves within 2-3 weeks, while the majority of visible swelling diminishes over the first 3-6 months.

Patients should expect to see a noticeable improvement in their nasal contour and, if applicable, nasal breathing within a few weeks, though the final refined aesthetic outcome will take longer to manifest as subtle swelling dissipates. The skin over the nose may feel numb initially, but sensation typically returns over several months. A successful course involves no significant complications such as infection, uncontrolled bleeding, or severe asymmetry, and the patient's satisfaction with both the functional and aesthetic improvements is high. Adherence to postoperative care instructions is paramount for achieving this expected positive outcome.

Comprehensive postoperative care and regular follow-up are critical for optimal healing and long-term results after rhinoplasty and septorhinoplasty.

\begin{itemize}
  \item **Initial Follow-up:** The first follow-up visit typically occurs around 5-7 days post-surgery for removal of the external splint, any external sutures, and sometimes internal splints. The surgeon will clean the nose and inspect the healing incisions.
  
  \item **Nasal Care:** Patients are instructed to use saline nasal sprays to keep the nasal passages moist and aid in crust removal. Gentle cleaning of the nostrils with cotton swabs may be advised.
  
  \item **Activity Restrictions:** Avoidance of strenuous physical activity, heavy lifting, and contact sports is crucial for at least 4-6 weeks to prevent trauma and bleeding. Nose-blowing should be avoided for several weeks.
  
  \item **Swelling Management:** Continued head elevation, especially while sleeping, and intermittent cold compresses (for the first few days) help reduce swelling. Taping of the nose may be recommended for several weeks to months to help control swelling and promote skin redraping.
  
  \item **Sun Protection:** The healing nasal skin is sensitive to sun exposure, so patients should use high-SPF sunscreen and avoid prolonged direct sunlight for at least 6 months to prevent hyperpigmentation.
  
  \item **Long-term Follow-up:** Subsequent follow-up visits are scheduled at 1 month, 3 months, 6 months, and 1 year post-surgery to monitor the resolution of swelling, assess the aesthetic and functional outcomes, and address any concerns.
  
  \item **Warning Signs:** Patients should contact their surgeon immediately if they experience severe or worsening pain, fever, purulent nasal discharge, uncontrolled bleeding, sudden vision changes, or signs of severe asymmetry or deformity.
\end{itemize}

\section*{Outcomes \& Prognosis}
Rhinoplasty consistently ranks among the top five most frequently performed cosmetic surgical procedures globally, with hundreds of thousands of operations performed annually. In the United States, it is a highly sought-after procedure, reflecting its significant impact on facial aesthetics and quality of life. The patient demographic is broad, but the majority of individuals undergoing rhinoplasty are young to middle-aged adults, typically between 20 and 40 years old. There is a slight female predominance for purely cosmetic rhinoplasty, though men frequently undergo the procedure for both aesthetic and functional reasons, particularly for correction of nasal obstruction.

Functional indications, such as septal deviation and nasal valve collapse, are exceedingly common, affecting a substantial portion of the general population to varying degrees. While the overall incidence of major complications is low, the rate of revision surgery for rhinoplasty is notably higher than for many other cosmetic procedures, ranging from approximately 5\% to 15\%. This reflects the inherent complexity of nasal anatomy, the subjective nature of aesthetic outcomes, and the challenges of predicting long-term tissue healing. Mortality associated with rhinoplasty is exceedingly rare, primarily linked to general anesthesia risks.

The overall outcomes and prognosis for rhinoplasty and septorhinoplasty are generally excellent, with high rates of patient satisfaction when realistic expectations are established preoperatively and the surgery is performed by an experienced, board-certified surgeon. Functional outcomes for nasal obstruction due to septal deviation or nasal valve collapse are highly favorable, with a significant majority of patients reporting substantial improvement in nasal breathing and quality of life. Aesthetic outcomes, while more subjective, are also largely positive, with patients achieving a more harmonious and balanced facial appearance.

The final results of rhinoplasty are typically stable for many years, though subtle changes can occur with aging. Key prognostic factors influencing the outcome include the surgeon's expertise, meticulous preoperative planning, the patient's skin type (thicker skin can mask subtle refinements and prolong swelling), cartilage quality and strength, and strict adherence to postoperative care instructions. While there isn't a single landmark clinical trial like NASCET for carotid endarterectomy, numerous studies and meta-analyses consistently demonstrate the efficacy and safety of modern rhinoplasty techniques. Despite high success rates, a small percentage of patients (5-15\%) may require revision surgery to address minor residual deformities, asymmetries, or persistent functional issues, underscoring the technical demands of the procedure and the dynamic nature of healing.

\section*{References}
\begin{enumerate}
  \item MedlinePlus. Rhinoplasty. U.S. National Library of Medicine; 2024. Available from: \url{https://medlineplus.gov/rhinoplasty.html}
  
  \item Rohrich RJ, Ahmad J. Rhinoplasty. Plast Reconstr Surg. 2011;128(2):49e-73e.
  
  \item Sabiston DC, Townsend CM. Sabiston Textbook of Surgery: The Biological Basis of Modern Surgical Practice. 21st ed. Elsevier; 2021.
  
  \item American Academy of Otolaryngology--Head and Neck Surgery. Position Statement: Rhinoplasty. 2019.
\end{enumerate}

\clearpage
